\label{sec:related_work}

\subsection{Model Predictive Control for Legged Robots}

Model predictive control (MPC) has become one of the dominant paradigms for dynamic legged locomotion because it provides a principled way to couple motion generation, contact scheduling, and constraint handling in a receding-horizon loop. Early influential work demonstrated that even simplified predictive models can produce highly effective locomotion behaviors, as exemplified by convex MPC for the MIT Cheetah 3 \cite{dicarlo2018convexmpc}. Since then, the field has progressively moved toward richer nonlinear and whole-body formulations. Grandia \emph{et al.} \cite{grandia2023perceptive} showed that nonlinear MPC can be combined with perceptive foothold reasoning to achieve dynamic locomotion over challenging terrain, while Meduri \emph{et al.} \cite{meduri2022biconmp} proposed a whole-body nonlinear MPC framework capable of generating both cyclic and acyclic motions online. In parallel, multicontact planning and whole-body optimization have been extended toward loco-manipulation settings \cite{sleiman2023versatile}, further increasing the dimensionality and multimodality of the underlying optimal control problem.

These works collectively establish that MPC is sufficiently expressive for agile legged behavior. However, the dominant emphasis in the literature has been on improving terrain awareness, whole-body planning, and task generality. Comparatively less attention has been paid to \emph{mechanism-level branch consistency}, i.e., whether the optimizer remains on the physically correct kinematic branch of a jointed mechanism throughout online optimization. Our work addresses this gap in the context of a hexapod, where the larger number of legs and contacts amplifies solver sensitivity to branch-inconsistent local minima.

\subsection{Local Minima in High-Dimensional Contact-Rich Optimization}

The difficulty of legged optimization does not arise solely from state dimension; it is also rooted in contact-induced nonconvexity and hybrid mode structure. Posa \emph{et al.} \cite{posa2014direct} highlighted that contact-rich trajectory optimization is inherently difficult because feasible motion depends on discontinuous interaction events, complementarity structure, and mode-dependent dynamics. More recent predictive-control formulations for agile maneuvers likewise acknowledge the challenge of solving highly nonlinear, multi-contact optimization problems online \cite{mastalli2022agile}. In such settings, local minima can be especially deceptive: two solutions may appear nearly equivalent in task space while corresponding to different---and not equally admissible---joint-space branches.

This issue is particularly problematic in warm-started SQP-style pipelines. When the local quadratic model is built around an already biased iterate, the solver can preserve a physically incorrect but locally attractive branch if that branch is not explicitly excised from the feasible set. Runtime guards may later intercept the resulting commands, but such post hoc filtering only hides the pathology rather than resolving it. Our results support the view that, for multi-legged systems, local-minimum sensitivity should be treated not merely as an optimization nuisance but as a \emph{safety-relevant modeling issue}.

\subsection{Interior-Point Methods and Hard Constraint Enforcement}

Interior-point methods provide an alternative perspective on constrained optimal control by embedding inequality satisfaction directly into the primal-dual solve rather than treating it as an external penalty preference. The modern nonlinear programming literature has long established the robustness of primal-dual interior-point methods with filter line search for large-scale constrained problems \cite{wachter2006ipopt}. In embedded and MPC-oriented settings, Frison and Diehl \cite{frison2020hpipm} showed that high-performance IPM implementations can solve structured optimal-control QPs efficiently while preserving strong constraint-treatment properties.

Despite these advantages, much of the legged-robot MPC literature has historically favored SQP, real-time iteration, or convex approximations because of their speed and implementation maturity. Our work suggests that this preference can become fragile when the controller must enforce \emph{mechanism-specific hard inequalities} whose violation leads to qualitatively wrong physical behavior. In contrast to prior work that uses hard constraints primarily for contact, friction, or foothold admissibility, we introduce a \emph{branch-consistency constraint} that encodes left/right tarsus sign admissibility directly in the optimal control problem. By pairing this constraint with an IPM backend, we ensure that branch validity is treated as a solver-level feasibility condition rather than as a runtime repair rule.

\subsection{Positioning of This Work}

The present paper differs from prior legged-MPC work in both diagnosis and remedy. On the diagnosis side, we identify a hidden failure mode in hexapod locomotion where the optimizer converges to a task-space-plausible but branch-inconsistent tarsus configuration. On the remedy side, we show that the combination of a solver-visible \emph{TarsusBranchConstraint} and an IPM backend eliminates this failure in real time. Therefore, our contribution is not simply another nonlinear MPC controller, but a feasible-set-centered reformulation of online legged optimization that connects physical branch validity, runtime safety, and solver design in a unified framework.
